\phantomsection\section*{\centering KẾT LUẬN}
\addcontentsline{toc}{section}{\numberline{}KẾT LUẬN}

\phantomsection\subsection*{Kết luận chung}
\addcontentsline{toc}{section}{\numberline{}Kết luận chung}

Đồ án đã hoàn thành mục tiêu trình bày và kiểm chứng một hướng mô phỏng kênh MIMO băng rộng dựa trên biểu diễn DPS. Mạch nội dung chính của báo cáo bắt đầu từ nút thắt tính toán của SoCE, sau đó khai thác đặc điểm giới hạn băng của kênh để xây dựng biểu diễn không gian con bằng DPS. Trên cơ sở đó, đồ án đã trình bày mô hình GCM MIMO băng rộng, quá trình rời rạc hóa bốn chiều và các nhánh mô phỏng MATLAB tương ứng.

Trong phạm vi kết quả đã kiểm chứng, SoCE được dùng làm kênh tham chiếu; DPS chính xác kiểm tra sai số cắt không gian con; DPS xấp xỉ 4D đánh giá ảnh hưởng của công thức xấp xỉ hệ số; và phương pháp hybrid kết hợp DPS theo thời gian--tần số với hàm mũ không gian trực tiếp. Các nhận xét định lượng trong báo cáo đều dựa trên các hình và bảng kết quả đã trình bày ở Chương 4.

Kết quả hiện tại cho thấy biểu diễn DPS có thể giảm số chiều biểu diễn khối kênh trong cấu hình mô phỏng đang xét, đồng thời phương pháp hybrid đạt sai số nhỏ so với SoCE tham chiếu. Tuy nhiên, kết luận này cần được hiểu trong phạm vi một cấu hình mô phỏng cụ thể, chưa phải là khẳng định tổng quát cho mọi kịch bản MIMO băng rộng.

\phantomsection\subsection*{Hướng phát triển}
\addcontentsline{toc}{section}{\numberline{}Hướng phát triển}

Các hướng phát triển chính gồm quét tham số số chiều DPS và hệ số phân giải, mở rộng kích thước mô phỏng, khớp lại một hình hoặc bảng cụ thể trong bài báo gốc, và chuẩn hóa mã MATLAB thành các hàm độc lập dễ tái sử dụng. Ngoài ra, có thể khảo sát thêm các phân bố MPC không đều hoặc các kịch bản kênh gần hơn với mô hình chuẩn để đánh giá tính ổn định của phương pháp.

\phantomsection\subsection*{Kiến nghị và đề xuất (nếu có)}
\addcontentsline{toc}{section}{\numberline{}Kiến nghị và đề xuất (nếu có)}

Trong các nghiên cứu tiếp theo, nên ưu tiên tái tạo một kết quả định lượng của bài báo Kaltenberger với bộ tham số được khớp đầy đủ trước khi mở rộng sang các cấu hình lớn hơn. Bước này giúp làm rõ mức độ tương thích giữa mô phỏng MATLAB của đồ án và kết quả gốc, từ đó tạo cơ sở chắc chắn hơn cho các so sánh về độ phức tạp và sai số.

\cleardoublepage
